\input{CSCSShortExample/cscsstyle}

\begin{document}

\Title{Efficent Simulation of Noise in Quantum Circuits}
\Author{Zolt\'an K\'egli, Tam\'as Kozsik, P\'eter Rakyta}

\begin{Abstract}
  Exact noisy quantum-circuit simulation remains expensive, yet many structured
  workloads contain small local noisy regions whose evolution may be reusable if
  their open-system semantics are preserved exactly. This paper studies a bounded
  form of that problem: same-support 1- and 2-qubit noisy motifs represented as
  composable CPTP objects and validated against sequential density-matrix
  evolution. The contribution is methodological rather than universal-performance
  driven. We show that such motifs can be promoted to exact reusable blocks,
  apply them on local support inside larger density states, and separate their
  evaluation into two layers: strict motif-level correctness and hybrid
  whole-workload continuity checks. The current evidence includes counted
  microcase validation, external-reference rows on the required slice, two
  whole-workload correctness anchors at 4 and 6 qubits, and a full emitted
  26-row workload matrix over frozen families. The resulting picture is mixed but
  scientifically useful: exact noisy fusion is feasible on the bounded support
  surface studied here, while workload-level acceleration remains unresolved and
  must be argued separately from correctness.
  \keywords{quantum circuit simulation, density-matrix simulation, noisy circuit fusion, CPTP channels, circuit partitioning, SQUANDER framework}
\end{Abstract}

\section{Introduction}

Partitioning and gate fusion are standard tools in unitary quantum simulation
\cite{r7,r8,r10,r9}, while open-system theory provides the language of Kraus maps,
Choi matrices, and superoperators for noisy evolution \cite{r2,r3}. What
remains underdeveloped is the
bridge between these two traditions: how to treat a small noisy region of a
circuit as an exact reusable object without losing the ordered semantics of
open-system evolution. The present study provides a bounded answer. It shows
that same-support 1- and 2-qubit noisy motifs can be composed as exact CPTP
objects and validated against sequential density-matrix evolution, rather than
being treated only as barriers between unitary islands. It also sharpens the
methodological picture: a scientifically credible path appears to need both a
strict motif-proof layer, where the fused object itself is validated, and a
separate hybrid whole-workload layer, where that object is tested inside larger
exact workloads with explicit route attribution. The current evidence now
reaches both layers: a bounded counted correctness package spans the strict
microcase surface plus two hybrid continuity anchors, the bounded
external-reference slice is present on the required rows, and a full emitted
26-row whole-workload matrix shows that feasibility does not automatically
imply performance benefit. The scientific value is therefore methodological and
decision-oriented: it identifies a reusable path toward exact noisy
acceleration while also showing, on a frozen workload family, where the
existing exact baseline is already sufficient and where the richer path is
still not justified.

\section{Motivation}

Exact noisy simulation matters because realistic local noise can qualitatively
change the behavior of variational algorithms. For questions about trainability,
loss landscapes, or optimizer robustness, a state-vector-only picture is often
not enough. Density-matrix simulation provides the cleanest exact reference, but
it is computationally expensive and therefore forces careful choices about what
kind of structure can be reused safely \cite{r4,r5,r6}.

This tension creates a natural scientific target: identify classes of noisy
structure that can be fused exactly, so that noisy simulation becomes more
structured than "unitaries plus interruptions" without sacrificing physical
correctness.

\section{Related Work}

The current literature offers strong ingredients, but not a fully integrated
answer to the bounded exact noisy-fusion question addressed here.

Graph-based partitioning and gate-fusion papers such as TDAG, GTQCP, QGo, and
QMin show how much can be gained when circuits are treated as structured
dependency objects \cite{r7,r8,r10,r9}. However, these methods are typically
developed for unitary or state-vector simulation.

Open-system and density-matrix literature, by contrast, gives the mathematical
language needed for exact noisy evolution. Nielsen and Chuang formalize quantum
channels and CPTP structure \cite{r2}; Wood, Biamonte, and Cory make explicit
the relations among Kraus, Choi, and Liouville forms \cite{r3}; Li et al.,
QuEST, and Qulacs show that exact mixed-state simulation is a serious
high-performance computing problem in its own right \cite{r4,r5,r6}.

What is still relatively unexplored is the exact noisy middle ground: bounded
fusion of mixed gate-noise motifs inside a partitioned execution flow.

\section{Methodology and result}

The current result is intentionally narrow but already meaningful.

It shows that a small same-support noisy motif can be promoted to a first-class
exact object: a bounded CPTP block composed in operation order and validated
against sequential density-matrix evolution. In the present bounded slice, this
has been demonstrated for 1-qubit and 2-qubit motifs and for local-support
application inside a larger density state.

Just as importantly, the current interpretation separates two scientific roles
that are often blurred together. One role is to prove that the exact noisy
object itself is sound. The other is to test whether that object helps inside
larger exact workloads that still contain structure outside the bounded fused
slice. The former belongs to a strict motif-proof setting; the latter belongs
to an explicit hybrid whole-workload setting.

This two-layer story is no longer only conceptual. The hybrid layer now has two
counted whole-workload correctness anchors at 4 and 6 qubits that match the
sequential oracle while preserving explicit route attribution between
channel-native and existing exact execution in SQUANDER \cite{r1}. The bounded
external-reference slice is also present on the required rows, which keeps the
current narrative tied to a real external check without turning the study into
a broad simulator comparison. In addition, the study now includes a full
emitted 26-row whole-workload matrix across the frozen workload families. That
matrix does not support a positive acceleration result. Instead, it shows a
bounded decision pattern: some workloads are already well served by the
existing exact fused baseline, while the richer path remains not yet justified
on the rest.

Scientifically, this matters because it changes the role of noise in the
simulation narrative. Noise is no longer treated only as a point where fusion
must stop. Instead, some noisy motifs can themselves become the fused object,
provided the representation, ordering, and invariant checks are all exact.

This is a result about \textbf{feasibility and methodology}, not about universal
performance. That distinction is important and should remain explicit.

\section{Discussion}

The current bounded study suggests five reusable principles for future exact
noisy acceleration work.

\begin{itemize}
  \item Start from the \textbf{ordered noisy semantics}, not from the representation alone.
        A fused object is only meaningful if it preserves the exact order of gates and
        channels.
  \item Choose one \textbf{primary exact representation} for the counted claim. Here the
        useful lesson is not that Kraus form is universally best, but that one primary
        representation should anchor the scientific claim.
  \item Require \textbf{physical invariant checks} before making performance claims.
        Trace-preservation and positivity are not optional bookkeeping; they are part
        of the evidence that the fused object remains a valid channel.
  \item Enforce \textbf{no silent fallback}. If a method advertises noisy fusion, then
        unsupported cases must remain visible. Otherwise benchmarks become difficult
        to interpret scientifically.
  \item Separate \textbf{proof mode} from \textbf{whole-workload evaluation mode}. The exact
        fused object can be scientifically established on a strict bounded slice
        before broader workload-level claims are made through an explicit hybrid
        interpretation.
\end{itemize}

These principles are potentially reusable beyond the present bounded support
surface. The full emitted matrix now reinforces that separation: exactness can
close while workload-level justification still remains negative or mixed.

Several hard questions remain open.

First, mathematical feasibility and performance usefulness are not the same.
Showing that a noisy motif can be fused exactly does not yet show that the
resulting method improves runtime or memory on the workloads that matter most.
The emitted whole-workload matrix makes that point concrete: the stronger
positive-method threshold is not met on the frozen slice.

Second, exact noisy acceleration remains scale-limited by dense mixed-state
representation itself. Even a successful bounded fusion method operates inside a
regime where memory and data movement remain fundamental constraints.

Third, broader noisy structure is still unresolved: correlated noise, larger
supports, and more heterogeneous motif families may require different
representations or different validation strategies.

\section{Conclusion}

The right scientific conclusion is therefore not "noisy fusion is solved," but
"a bounded exact path now exists and can be studied honestly."

\begin{thebibliography}{99}

  \bibitem{r1}
  SQUANDER: \texttt{https://github.com/rakytap/sequential-quantum-gate-decomposer}

  \bibitem{r2}
  Michael A. Nielsen and Isaac L. Chuang, \textit{Quantum Computation and Quantum
    Information}, Cambridge University Press (2010).

  \bibitem{r3}
  Christopher J. Wood, Jacob D. Biamonte, and David G. Cory, \textit{Tensor networks
    and graphical calculus for open quantum systems}, \texttt{Quantum Information and
    Computation 15, 759-811 (2015)}.

  \bibitem{r4}
  Ang Li, Omer Subasi, Xiu Yang, and Sriram Krishnamoorthy, \textit{Density Matrix
    Quantum Circuit Simulation via the BSP Machine on Modern GPU Clusters}, SC20.

  \bibitem{r5}
  Tyson Jones, Anna Brown, Ian Bush, and Simon C. Benjamin, \textit{QuEST and High
    Performance Simulation of Quantum Computers}, \texttt{Scientific Reports 9, 10736
    (2019)}.

  \bibitem{r6}
  Yasunari Suzuki et al., \textit{Qulacs: a fast and versatile quantum circuit
    simulator for research purpose}, \texttt{Quantum 5, 559 (2021)}.

  \bibitem{r7}
  Joseph Clark, Travis S. Humble, and Himanshu Thapliyal, \textit{TDAG: Tree-based
    Directed Acyclic Graph Partitioning for Quantum Circuits}, ACM GLSVLSI 2023.

  \bibitem{r8}
  Joseph Clark, Travis S. Humble, and Himanshu Thapliyal, \textit{GTQCP: Greedy
    Topology-Aware Quantum Circuit Partitioning}, \texttt{arXiv:2410.02901}.

  \bibitem{r10}
  Xin-Chuan Wu, Marc Grau Davis, Frederic T. Chong, and Costin Iancu, \textit{QGo:
    Scalable Quantum Circuit Optimization Using Automated Synthesis},
  \texttt{arXiv:2012.09835}.

  \bibitem{r9}
  Longshan Xu, Edwin Hsing-Mean Sha, Yuhong Song, and Qingfeng Zhu, \textit{QMin:
    Quantum Circuit Minimization via Gate Fusions for Efficient State Vector
    Simulation}, \texttt{Quantum Information Processing 25, 6 (2026)}.

\end{thebibliography}
\end{document}
